\documentclass[11pt]{article}

\usepackage{amsmath}
\usepackage{amssymb}
\usepackage{amsthm}
\usepackage[hidelinks]{hyperref}

\newtheorem{theorem}{Theorem}[section]
\newtheorem{lemma}[theorem]{Lemma}
\newtheorem{proposition}[theorem]{Proposition}
\newtheorem{corollary}[theorem]{Corollary}
\newtheorem{conjecture}[theorem]{Conjecture}

\theoremstyle{definition}
\newtheorem{hypothesis}[theorem]{Hypothesis}
\newtheorem{definition}[theorem]{Definition}

\theoremstyle{remark}
\newtheorem{remark}[theorem]{Remark}
\newtheorem{example}[theorem]{Example}

\newcommand{\Z}{\mathbb{Z}}
\newcommand{\Q}{\mathbb{Q}}
\newcommand{\F}{\mathbb{F}}
\newcommand{\legendre}[2]{\left(\tfrac{#1}{#2}\right)}
\DeclareMathOperator{\Res}{Res}

\title{A class number criterion for pentads of congruent factorials}
\author{Jeff Hinojosa\thanks{With machine-assisted proof search and verification.
All computations reported here, and substantial portions of the proof search and
adversarial checking, were carried out by AI agents (Claude, Anthropic) under the
author's direction; this disclosure follows current norms for AI assistance in
mathematical work. See Remark~\ref{rem:provenance} for details on verification
status.}}
\date{July 23, 2026}

\begin{document}

\maketitle

\begin{abstract}
Let $q \equiv 1 \pmod 6$, $q \ge 7$, and let $p \equiv 3 \pmod 4$ be a prime
divisor of $q!-1$; such a divisor always exists, and any prime divisor of
$q!-1$ exceeds $q$. Building on the tetrad construction of
Agust\'in-Aquino and Hern\'andez Santiago, we show that the four residues
$1$, $q$, $p-1-q$, $p-2$ are pairwise distinct with all factorials
$\equiv 1 \pmod p$, and that the central residue $m = (p-1)/2$ joins them ---
producing five equal factorials, hence $k = 4$ consecutive intervals with
product $\equiv 1 \pmod p$ in Erd\H{o}s problem \#1056 --- \emph{if and only
if} the class number of $\Q(\sqrt{-p})$ satisfies $h(-p) \equiv 3 \pmod 4$.
The proof is elementary apart from the counting form of Dirichlet's class
number formula, and it recovers a theorem of Mordell. We deduce a conditional
infinitude theorem for pentad primes, a multi-seed shape theorem that reduces
hexad infinitude to the statement that infinitely many pairs
$q_1 < q_2 \equiv 1 \pmod 6$ satisfy $\gcd(q_1!-1,\,q_2!-1) > 1$, and a
limitative theorem explaining why the classical identity toolkit terminates,
unconditionally, exactly at tetrads. Twelve pentads are certified end to end;
the smallest is $p = 359$ and the most striking is $p = 5039 = 7!-1$, whose
complete solution set of $n! \equiv 1 \pmod{5039}$ is precisely the
constructed pentad. The unconditional infinitude question for $k \ge 4$
remains open.
\end{abstract}

\section{Introduction}\label{sec:intro}

Erd\H{o}s asked (problem \#1056 in the collection at
\href{https://www.erdosproblems.com/1056}{erdosproblems.com}; problem A15 in
Guy \cite{Guy}): for $k \ge 2$, is there a prime $p$ and $k$ consecutive
blocks of consecutive integers
\[
  I_1 = (q_0, q_1], \quad I_2 = (q_1, q_2], \quad \dots, \quad
  I_k = (q_{k-1}, q_k], \qquad 1 \le q_0 < q_1 < \dots < q_k \le p-1,
\]
such that the product of the elements of each block is $\equiv 1 \pmod p$?
Since $\prod_{n \in (q_{i-1}, q_i]} n = q_i!/q_{i-1}!$, the existence of $k$
such blocks is equivalent to the existence of $k+1$ values
$q_0 < q_1 < \dots < q_k$ in $[1, p-1]$ whose factorials are congruent to a
common residue modulo $p$. We call a set of $j$ such values with equal
factorials a $j$-\emph{ad}: tetrads ($j=4$) give $k=3$, pentads ($j=5$) give
$k=4$, hexads ($j=6$) give $k=5$.

Witnesses are known computationally for all $k \le 14$: OEIS sequence
A060427 \cite{OEIS} records the least prime for each $k$, beginning
\[
  11,\ 17,\ 23,\ 71,\ 71,\ 599,\ 599,\ 3011,\ 27901,\ 52163,\ 778699,\
  2374649,\ 10428007
\]
for $k = 2, \dots, 14$ (OEIS \texttt{A060427}; extended by
J.~K.~Andersen, 2007). Whether witnesses exist for
\emph{every} $k$, and whether any single $k \ge 2$ admits infinitely many
primes, are the substantive open questions. In 2026, Agust\'in-Aquino and
Hern\'andez Santiago \cite{AAHS} proved, with a Lean-verified development,
that infinitely many primes carry a tetrad: for $q \equiv 1 \pmod 6$,
$q \ge 7$, and any prime $p \mid q!-1$, the four residues
$\{1, q, p-1-q, p-2\}$ all have factorial $\equiv 1 \pmod p$ and are pairwise
distinct. This was the state of the art: $k = 3$ with infinitely many
witnesses, and nothing beyond.

This note adds one point to their configuration and finds that the question
of whether it joins is governed by a class number. Our main result is the
following.

\begin{theorem}[Pentad criterion; proved as Theorem~\ref{thm:main}]
Let $q \equiv 1 \pmod 6$ with $q \ge 7$, and let $p \equiv 3 \pmod 4$ be a
prime divisor of $q!-1$. (Such a divisor always exists, since
$q!-1 \equiv 3 \pmod 4$; and necessarily $p > q$.) Put $m = (p-1)/2$. Then
the five residues
\[
  \{\,1,\ q,\ m,\ p-1-q,\ p-2\,\} \subset [1, p-1]
\]
are pairwise distinct, the four outer values satisfy
$x! \equiv 1 \pmod p$, and $m! \equiv \pm 1 \pmod p$, with
\[
  m! \equiv 1 \pmod p
  \iff
  h(-p) \equiv 3 \pmod 4,
\]
where $h(-p)$ is the class number of $\Q(\sqrt{-p})$. In that case the five
values form a pentad and $p$ carries $k = 4$ consecutive intervals of product
$\equiv 1 \pmod p$; otherwise $m! \equiv -1 \pmod p$ and $p$ carries the
tetrad together with the pair $\{m, p-1\}$ at common value $-1$.
\end{theorem}

The mechanism is a theorem of Mordell \cite{Mordell}: for $p \equiv 3 \pmod
4$, $\left(\frac{p-1}{2}\right)! \equiv (-1)^{(h(-p)+1)/2} \pmod p$. We
re-prove it here in an elementary two-step form --- first
$m! \equiv (-1)^{N_p} \pmod p$ where $N_p$ counts quadratic nonresidues in
$(0, p/2)$, then $N_p$ even $\iff h(-p) \equiv 3 \pmod 4$ via the counting
form of Dirichlet's class number formula --- because both halves are needed
separately in what follows. The point of the theorem is not the central-sign
law, which is classical, but the coupling: \emph{Erd\H{o}s \#1056 at $k = 4$
is tied, along an explicit infinite family of primes, to the residue of
$h(-p)$ modulo $4$} --- a quantity about which, beyond Gauss's parity result
($h(-p)$ is odd for prime $p \equiv 3 \bmod 4$), no unconditional
determination of $h(-p) \bmod 4$ is known.
Indeed, whether \emph{each} residue class $h(-p) \equiv 1, 3 \pmod 4$
occurs for infinitely many primes $p \equiv 3 \pmod 4$ appears to be open
even without any restriction to a thin family (by pigeonhole at least one
class is infinite, but it is not known which); the question circulates in
this form on MathOverflow since 2013, unresolved. A structural reason is
supplied by Williams \cite{Wi79}, who proved
$h(-p) \equiv h(p) + U + 1 \pmod 4$, where $T + U\sqrt{p}$ is the
fundamental unit of $\mathbb{Q}(\sqrt{p})$: the invariant is governed by
Pell-equation and real-quadratic data rather than by any splitting
condition on $p$, so no Chebotarev-type criterion is available. In this
sense the hypothesis of our conditional infinitude theorem
(Section~\ref{sec:corollaries}) is optimal: it isolates precisely the
recognized open difficulty, with nothing else in the way.

The note is organized as follows. Section~\ref{sec:main} states the main
theorem and fixes notation; Section~\ref{sec:proofs} gives complete proofs.
Section~\ref{sec:corollaries} derives the conditional infinitude of pentad
primes, an unconditional dichotomy, the multi-seed shape theorem (every
$k \ge 3$ is realized by some seed pattern), and the reduction of hexad
infinitude to a clean divisibility hypothesis with numerical witnesses.
Section~\ref{sec:limits} proves a limitative theorem: the identity method
that produces tetrads has a built-in ceiling and cannot self-start beyond
them, so that pentad or hexad infinitude \emph{via identities at the pinned
value $1$} is equivalent to the class number statement or to multi-seed
divisibilities. The ceiling is a feature of the pinned value: \#1056 permits
a free common residue, and Remark~\ref{rem:freevalue} exhibits $k = 5$
configurations off both $S_1(p)$ and $S_{-1}(p)$, where the analysis does
not reach.
Section~\ref{sec:numerics} reports the numerical evidence: twelve certified
pentads, the flagship $p = 5039 = 7!-1$, and sign statistics.
Remark~\ref{rem:provenance} states plainly what has been machine-verified and
what has not. We emphasize at the outset: \emph{the tetrad construction is
entirely due to Agust\'in-Aquino and Hern\'andez Santiago \cite{AAHS}; our
contribution is the central point, the class number dichotomy, and the
limitative analysis. The unconditional infinitude of witnesses for any
$k \ge 4$ remains open.}

\section{The main theorem}\label{sec:main}

\subsection*{Notation}
Throughout, $p$ denotes an odd prime and $m = (p-1)/2$. For a residue $r$ we
write $S_r(p) = \{\, n \in [1, p-1] : n! \equiv r \pmod p \,\}$. The Legendre
symbol is $\legendre{a}{p}$, and
\[
  N_p = \#\{\, a : 0 < a < p/2,\ \legendre{a}{p} = -1 \,\}, \qquad
  A_p = \#\{\, a : 0 < a < p/2,\ \legendre{a}{p} = +1 \,\},
\]
so $A_p + N_p = m$. We write $h(-p)$ for the class number of the imaginary
quadratic field $\Q(\sqrt{-p})$ (equivalently, of the discriminant $-p$ when
$p \equiv 3 \bmod 4$). A \emph{seed} for $p$ is an integer $q$ with
$p \mid q! - 1$.

\begin{theorem}[Pentad criterion]\label{thm:main}
Let $q \equiv 1 \pmod 6$, $q \ge 7$, and let $p \equiv 3 \pmod 4$ be a prime
with $p \mid q! - 1$. Then:
\begin{enumerate}
  \item[(a)] such a prime $p$ exists for every such $q$, and every prime
    divisor of $q!-1$ satisfies $p > q$ (indeed $p \ge q+4$);
  \item[(b)] the five integers $1,\ q,\ m,\ p-1-q,\ p-2$ are pairwise
    distinct elements of $[1, p-1]$;
  \item[(c)] $1! \equiv q! \equiv (p-1-q)! \equiv (p-2)! \equiv 1 \pmod p$,
    and $m! \equiv \pm 1 \pmod p$;
  \item[(d)] $m! \equiv 1 \pmod p \iff N_p$ is even $\iff h(-p) \equiv 3
    \pmod 4$.
\end{enumerate}
Consequently, exactly one of the following holds:
\begin{enumerate}
  \item[(i)] $h(-p) \equiv 3 \pmod 4$: then $\{1, q, m, p-1-q, p-2\}$,
    sorted increasingly, is a pentad at common value $1$, and $p$ carries
    $k = 4$ consecutive intervals each of product $\equiv 1 \pmod p$;
  \item[(ii)] $h(-p) \equiv 1 \pmod 4$: then $m! \equiv -1 \pmod p$, and $p$
    carries the tetrad at $+1$ together with the pair $\{m, p-1\}$ at $-1$.
\end{enumerate}
\end{theorem}

Part (a) and the tetrad portion of (b)--(c) constitute the theorem of
\cite{AAHS}, refined by the observation --- available at zero extra cost ---
that the prime $p$ may always be chosen in the class $3 \bmod 4$; the
statement of \cite{AAHS} imposes no congruence condition on $p$. Parts (b)
and (d) for the central point $m$, and the resulting dichotomy, are new to
the best of our knowledge (see Remark~\ref{rem:provenance} on the scope of
the novelty check).

\section{Proofs}\label{sec:proofs}

\subsection{The reflection lemma}

The engine of every construction in this subject is the Wilson pairing
identity. It is classical; we state it with the consequences we use.

\begin{lemma}[Reflection]\label{lem:reflection}
Let $p$ be an odd prime and $0 \le n \le p-1$. Then
\[
  n! \, (p-1-n)! \equiv (-1)^{n+1} \pmod p.
\]
Consequently:
\begin{enumerate}
  \item[(i)] if $n$ is odd and $n! \equiv c \pmod p$, then
    $(p-1-n)! \equiv c^{-1} \pmod p$; in particular the sets $S_1(p)$ and
    $S_{-1}(p)$ are stable under $n \mapsto p-1-n$ restricted to odd $n$;
  \item[(ii)] if $n$ is even and $n! \equiv 1$, then $(p-1-n)! \equiv -1$
    (and vice versa);
  \item[(iii)] $1, \, p-2 \in S_1(p)$ and $p-1 \in S_{-1}(p)$ for $p \ge 5$;
  \item[(iv)] $(p-k)! \equiv (-1)^k \, \bigl((k-1)!\bigr)^{-1} \pmod p$ for
    $1 \le k \le p-1$; hence $(p-k)! \equiv 1$ is equivalent to
    $p \mid (k-1)! - (-1)^k$.
\end{enumerate}
\end{lemma}

\begin{proof}
Write $(p-1)! = n! \prod_{j=n+1}^{p-1} j$. For $n+1 \le j \le p-1$ put
$i = p - j$, so $1 \le i \le p-1-n$ and $j \equiv -i \pmod p$; hence
\[
  \prod_{j=n+1}^{p-1} j \;\equiv\; (-1)^{p-1-n} (p-1-n)! \;\equiv\;
  (-1)^{n} (p-1-n)! \pmod p,
\]
since $p - 1$ is even. Wilson's theorem $(p-1)! \equiv -1$ gives the
identity. (i): $n$ odd makes the right side $+1$, so
$(p-1-n)! \equiv (n!)^{-1} = c^{-1}$; since $c = \pm 1$ are the only
self-inverse residues in $\F_p$, exactly the classes $S_{\pm 1}$ are
preserved (the class of $c \ne \pm 1$ is carried to the class of $c^{-1} \ne
c$). (ii): $n$ even makes the right side $-1$, so $(p-1-n)! \equiv -c^{-1}$.
(iii): $1! = 1$; from $(p-1)! = (p-1)(p-2)! \equiv -(p-2)! \equiv -1$ we get
$(p-2)! \equiv 1$; and $(p-1)! \equiv -1$ is Wilson.
(iv): $(p-1)! = (p-1)(p-2) \cdots (p-(k-1)) \cdot (p-k)! \equiv
(-1)^{k-1}(k-1)! \, (p-k)! \equiv -1$, so
$(p-k)! \equiv (-1)^k ((k-1)!)^{-1}$.
\end{proof}

\subsection{Existence of the seed prime, and its size}

\begin{lemma}[Seed prime]\label{lem:seed}
For every integer $q \ge 4$, the number $q! - 1$ satisfies
$q! - 1 \equiv 3 \pmod 4$, and hence has at least one prime factor
$p \equiv 3 \pmod 4$. Moreover every prime factor of $q!-1$ exceeds $q$.
\end{lemma}

\begin{proof}
Since $q \ge 4$, $4 \mid q!$, so $q! - 1 \equiv 3 \pmod 4$. An odd integer
all of whose prime factors are $\equiv 1 \pmod 4$ is itself $\equiv 1 \pmod
4$, because the class of $1$ is closed under multiplication in
$(\Z/4\Z)^{\times}$; hence $q!-1$ has a prime factor $p \equiv 3 \pmod 4$.
If a prime $p \le q$ divided $q!-1$, it would divide $q!$ and hence their
difference $1$, which is absurd; so $p > q$.
\end{proof}

\begin{remark}\label{rem:mod8}
The lemma cannot be sharpened to a class modulo $8$: indeed
$q! - 1 \equiv 7 \pmod 8$ for $q \ge 4$, but $3 \cdot 5 \equiv 7 \pmod 8$
shows that an integer $\equiv 7 \pmod 8$ need not have any prime factor
$\equiv 7 \pmod 8$. The modulo-$4$ conclusion is exactly what
Mordell's theorem needs, and exactly what is available.
\end{remark}

\subsection{Distinctness}

\begin{lemma}[Distinctness]\label{lem:distinct}
Let $q \equiv 1 \pmod 6$, $q \ge 7$, and let $p > q$ be an odd prime,
$m = (p-1)/2$. Then the five integers $1, q, m, p-1-q, p-2$ are pairwise
distinct elements of $[1, p-1]$, and moreover $p \ge q + 4$.
\end{lemma}

\begin{proof}
The pairwise coincidences among the four tetrad values reduce to six
equations:
\[
  1 = q; \qquad
  1 = p-1-q \iff p = q+2; \qquad
  1 = p-2 \iff p = 3;
\]
\[
  q = p-1-q \iff p = 2q+1; \qquad
  q = p-2 \iff p = q+2; \qquad
  p-1-q = p-2 \iff q = 1.
\]
Adjoining $m$ adds four more:
\[
  m = 1 \iff p = 3; \qquad
  m = p-2 \iff p = 3; \qquad
  m = q \iff p = 2q+1; \qquad
  m = p-1-q \iff p = 2q+1.
\]
Now $q \equiv 1 \pmod 6$ forces $q + 2 \equiv 3 \pmod 6$ and
$2q + 1 \equiv 3 \pmod 6$; both exceed $3$, so both are divisible by $3$ and
composite --- neither can equal the prime $p$. Also $p > q \ge 7 > 3$
excludes $p = 3$, and $q = 1$ is excluded by hypothesis. Hence no coincidence
occurs. Finally, $p$ is odd and $p \notin \{q+1, q+3\}$ (both even) and
$p \ne q + 2$, so $p \ge q + 4$; in particular $p - 1 - q \ge 3 \ge 1$ and
all five values lie in $[1, p-1]$.
\end{proof}

\subsection{The central sign law}

\begin{lemma}[Central sign]\label{lem:centralsign}
Let $p \equiv 3 \pmod 4$ be prime, $p > 3$, and $m = (p-1)/2$ (which is odd).
Then
\[
  m! \equiv (-1)^{N_p} \pmod p .
\]
\end{lemma}

\begin{proof}
Let $R \subset \F_p^\times$ be the set of quadratic residues, a subgroup of
odd order $|R| = m$. In a finite abelian group of odd order the product of
all elements is the identity: elements pair off with their inverses, and no
element other than the identity is self-inverse. Hence
$\prod_{r \in R} r \equiv 1 \pmod p$.

Since $p \equiv 3 \pmod 4$ we have $\legendre{-1}{p} = -1$, so the map
$a \mapsto p - a$ exchanges residues and nonresidues; in particular it maps
the nonresidues in $(0, p/2)$ bijectively onto the residues in $(p/2, p)$.
Therefore
\[
  R \;=\; \bigl(R \cap (0, p/2)\bigr) \ \sqcup \
  \bigl\{\, p - a : 0 < a < p/2, \ \legendre{a}{p} = -1 \,\bigr\},
\]
and consequently
\[
  1 \;\equiv\; \prod_{r \in R} r
  \;\equiv\;
  \Bigl(\prod_{\substack{0 < a < p/2 \\ \legendre{a}{p} = 1}} a\Bigr)
  \prod_{\substack{0 < a < p/2 \\ \legendre{a}{p} = -1}} (p - a)
  \;\equiv\;
  (-1)^{N_p} \prod_{0 < a < p/2} a
  \;=\; (-1)^{N_p} \, m! \pmod p .
\]
Since $(-1)^{N_p}$ is its own inverse, $m! \equiv (-1)^{N_p} \pmod p$.
(Consistently, Lemma~\ref{lem:reflection} at the odd fixed point $n = m$
gives $(m!)^2 \equiv 1$.)
\end{proof}

\subsection{The class number bridge}

\begin{lemma}[Class number bridge]\label{lem:bridge}
Let $p \equiv 3 \pmod 4$ be prime, $p > 3$. Then $h(-p)$ is odd, and
\[
  N_p \text{ is even} \iff h(-p) \equiv 3 \pmod 4 .
\]
Combined with Lemma~\ref{lem:centralsign} this gives Mordell's theorem
\emph{\cite{Mordell}}:
\[
  m! \;\equiv\; (-1)^{(h(-p)+1)/2} \pmod p ,
  \qquad\text{i.e.}\qquad
  m! \equiv +1 \iff h(-p) \equiv 3 \pmod 4 .
\]
\end{lemma}

\begin{proof}
Oddness of $h(-p)$ is Gauss's genus theory: the discriminant $-p$ has a
single prime discriminant divisor, so there is one genus, and the number of
classes per genus is odd; equivalently, the $2$-part of the class group is
trivial.

For the mod-$4$ statement we use the counting form of Dirichlet's class
number formula for a prime discriminant $-p < -4$ (see Landau
\cite{Landau} or Borevich--Shafarevich \cite{BS}):
\[
  h(-p) \;=\; \frac{A_p - N_p}{\,2 - \legendre{2}{p}\,} .
\]
Recall $A_p + N_p = m$. There are two branches.

\emph{Case $p \equiv 7 \pmod 8$.} Then $\legendre{2}{p} = +1$ and
$h = A_p - N_p = m - 2N_p$. Writing $p = 8s + 7$ gives $m = 4s + 3 \equiv 3
\pmod 4$, so $h \equiv 3 - 2N_p \pmod 4$: $h \equiv 3 \pmod 4$ iff
$2N_p \equiv 0 \pmod 4$ iff $N_p$ is even.

\emph{Case $p \equiv 3 \pmod 8$.} Then $\legendre{2}{p} = -1$ and
$3h = m - 2N_p$. Writing $p = 8s + 3$ gives $m = 4s + 1 \equiv 1 \pmod 4$,
so $3h \equiv 1 - 2N_p \pmod 4$ and, multiplying by $3^{-1} \equiv 3
\pmod 4$,
\[
  h \;\equiv\; 3(1 - 2N_p) \;\equiv\; 3 + 2N_p \pmod 4 :
\]
again $h \equiv 3 \pmod 4$ iff $N_p$ is even.

In both branches $N_p$ even $\iff h(-p) \equiv 3 \pmod 4$. Since $h$ is odd,
$(-1)^{(h+1)/2} = +1 \iff h \equiv 3 \pmod 4$, which matches
$m! \equiv (-1)^{N_p}$ from Lemma~\ref{lem:centralsign}.
\end{proof}

\begin{remark}
In our numerical verification (Section~\ref{sec:numerics}) the class numbers
were computed by exact counting of reduced primitive binary quadratic forms
of discriminant $-p$, independently cross-checked against the Dirichlet
character sum; the verification of the criterion therefore does not assume
the analytic formula it uses.
\end{remark}

\subsection{Proof of Theorem~\ref{thm:main}}

\begin{proof}[Proof of Theorem~\ref{thm:main}]
(a) is Lemma~\ref{lem:seed} (with $p \ge q+4$ from
Lemma~\ref{lem:distinct}); (b) is Lemma~\ref{lem:distinct}.

(c): $1! = 1$; $q! \equiv 1$ because $p \mid q! - 1$; $q$ is odd, so
Lemma~\ref{lem:reflection}(i) gives $(p-1-q)! \equiv 1$; and
$(p-2)! \equiv 1$ by Lemma~\ref{lem:reflection}(iii). Since
$p \equiv 3 \pmod 4$, $m$ is odd, and Lemma~\ref{lem:reflection} at $n = m$
gives $(m!)^2 \equiv (-1)^{m+1} = 1$, so $m! \equiv \pm 1 \pmod p$.

(d) is Lemma~\ref{lem:centralsign} combined with Lemma~\ref{lem:bridge}.

For the dichotomy: in case (i) the five distinct values, sorted as
$q_0 < q_1 < \dots < q_4$, all lie in $S_1(p)$, so each consecutive interval
$(q_{i-1}, q_i]$ has product $q_i!/q_{i-1}! \equiv 1 \pmod p$ --- four
intervals, i.e.\ $k = 4$ in the original formulation. In case (ii),
$m! \equiv -1$ places $m$ with the Wilson point $p-1$ in $S_{-1}(p)$. The two
cases are exhaustive and exclusive by (c) and (d).
\end{proof}

\section{Corollaries: conditional infinitude, shapes, and hexads}
\label{sec:corollaries}

\subsection{Pentad primes, conditionally}

Let
\[
  \mathcal F \;=\; \bigl\{\, (q, p) \ :\ q \equiv 1 \ (\mathrm{mod}\ 6),\
  q \ge 7,\ p \ \text{prime},\ p \mid q! - 1,\ p \equiv 3 \ (\mathrm{mod}\
  4) \,\bigr\}.
\]
By Lemma~\ref{lem:seed}, $\mathcal F$ projects onto all admissible $q$, and
$p > q$ throughout.

\begin{hypothesis}[CFS: central factorial sign]\label{hyp:cfs}
$h(-p) \equiv 3 \pmod 4$ for infinitely many pairs $(q, p) \in \mathcal F$.
\end{hypothesis}

Full equidistribution --- each residue $1, 3 \bmod 4$ with density $1/2$ ---
is what randomness models for class numbers predict, and what the data of
Section~\ref{sec:numerics} show; but only ``infinitely often'' is needed.

\begin{corollary}[Conditional infinitude of pentad primes]\label{cor:cond}
Assume Hypothesis~\ref{hyp:cfs}. Then infinitely many primes $p$ admit five
distinct values $x \in [1, p-1]$ with $x! \equiv 1 \pmod p$; that is,
Erd\H{o}s \#1056 holds for $k = 4$ with infinitely many prime witnesses.
\end{corollary}

\begin{proof}
For each pair $(q, p) \in \mathcal F$ with $h(-p) \equiv 3 \pmod 4$,
Theorem~\ref{thm:main}(i) produces a pentad for $p$. It remains to see the
primes involved are unbounded. For fixed $q$ the number $q!-1$ has finitely
many prime factors, so an infinite set of pairs contains pairs with $q$
arbitrarily large; and $p > q$ forces the corresponding primes to be
unbounded, hence infinitely many distinct primes occur.
\end{proof}

\begin{corollary}[Unconditional dichotomy]\label{cor:dichotomy}
At least one of the following is true:
\begin{enumerate}
  \item[(i)] infinitely many primes admit pentads at value $1$ ($k = 4$
    with infinitely many witnesses); or
  \item[(ii)] $h(-p) \equiv 1 \pmod 4$ for all but finitely many pairs
    $(q, p) \in \mathcal F$ --- an eventually-constant law for the residue
    of an odd class number along an explicit unbounded family of primes,
    contrary to every randomness model for $h(-p)$.
\end{enumerate}
\end{corollary}

\begin{proof}
If case (i) of Theorem~\ref{thm:main} occurs for infinitely many members of
$\mathcal F$, the argument of Corollary~\ref{cor:cond} gives conclusion (i);
if it occurs for only finitely many, conclusion (ii) holds by definition.
\end{proof}

\subsection{The multi-seed shape theorem}

The pentad construction generalizes: each additional seed buys more points,
with all degeneracies killed automatically by the mod-$6$ arithmetic.

\begin{theorem}[Shape theorem]\label{thm:shape}
Let $j \ge 1$ and $e \ge 0$. Choose distinct integers
$q_1 < \dots < q_j$, each $\equiv 1 \pmod 6$ and $\ge 7$ (hence odd), and
distinct even integers $t_1, \dots, t_e \ge 4$ (no congruence condition
needed). Suppose a prime $p$ divides $q_i! - 1$ for every
$i \in \{1, \dots, j\}$ and $t_\ell! - 1$ for every
$\ell \in \{1, \dots, e\}$. Then $p$ exceeds every seed, and the $2j + e + 2$
integers
\[
  \{1,\ p-2\} \ \cup\ \{q_1, \dots, q_j\} \ \cup\ \{t_1, \dots, t_e\}
  \ \cup\ \{p-1-q_1, \dots, p-1-q_j\}
\]
are pairwise distinct elements of $[1, p-1]$ with all factorials
$\equiv 1 \pmod p$. Sorted increasingly, they yield $k = 2j + e + 1$
consecutive intervals each of product $\equiv 1 \pmod p$.

\emph{Value-$(-1)$ variant:} if instead $p \mid q_i! + 1$ and
$p \mid t_i! + 1$, then $\{p-1\} \cup \{q_i\} \cup \{t_i\} \cup
\{p-1-q_i\}$ are $2j + e + 1$ distinct points with all factorials
$\equiv -1$, yielding $k = 2j + e$ intervals.

In particular $(j, e) = (1, 0)$ recovers the tetrad theorem of \cite{AAHS}
($k = 3$); $(1, 1)$ gives pentads ($k = 4$); $(2, 0)$ gives hexads
($k = 5$); and every $k \ge 3$ (indeed every $k \ge 2$, using the $-1$
variant) is realized by some $(j, e)$.
\end{theorem}

\begin{proof}
\emph{Size of $p$:} if $p \le q$ for some seed $q$, then $p \mid q!$ and
$p \mid q! \mp 1$ force $p \mid 1$, absurd; so $p$ exceeds every seed, and
since $p$ is odd, $p \ge (\text{largest odd seed}) + 2$, so all listed points
lie in $[1, p-1]$.

\emph{Factorial values} ($+1$ case): $1! = 1$; $q_i! \equiv 1$ and
$t_i! \equiv 1$ by hypothesis; $(p-2)! \equiv 1$ by
Lemma~\ref{lem:reflection}(iii); $(p-1-q_i)! \equiv 1$ by
Lemma~\ref{lem:reflection}(i), since $q_i$ is odd.

\emph{Distinctness} --- the complete case list:
\begin{enumerate}
  \item[{[a]}] $1 = q_i$ or $1 = t_i$: impossible, all seeds are $\ge 4$.
  \item[{[b]}] $1 = p-2 \iff p = 3$: impossible, $p$ exceeds the seeds.
  \item[{[c]}] $1 = p-1-q_i \iff p = q_i + 2$: but $q_i \equiv 1 \pmod 6$
    makes $q_i + 2 \equiv 3 \pmod 6$, so $3 \mid p$, forcing $p = 3$:
    impossible.
  \item[{[d]}] $q_i = q_{i'}$ or $t_i = t_{i'}$ ($i \ne i'$): excluded by
    choice.
  \item[{[e]}] $q_i = t_{i'}$: opposite parities.
  \item[{[f]}] $q_i = p-2 \iff p = q_i + 2$: as [c].
  \item[{[g]}] $t_i = p-2 \iff p = t_i + 2$, which is even: impossible for
    odd $p$.
  \item[{[h]}] $q_i = p-1-q_{i'} \iff p = q_i + q_{i'} + 1$: for $i = i'$
    this is $p = 2q_i + 1 \equiv 3 \pmod 6$; for $i \ne i'$ it is
    $\equiv 1 + 1 + 1 \equiv 3 \pmod 6$; either way $3 \mid p$: impossible.
  \item[{[i]}] $t_i = p-1-q_{i'} \iff p = q_{i'} + t_i + 1 =
    \text{odd} + \text{even} + 1$, which is even: impossible.
  \item[{[j]}] $p-1-q_i = p-1-q_{i'} \iff q_i = q_{i'}$.
  \item[{[k]}] $p-1-q_i = p-2 \iff q_i = 1$: impossible.
\end{enumerate}
So all $2j+e+2$ points are distinct, and consecutive interval products are
ratios of equal factorials, hence $\equiv 1$.

\emph{$(-1)$ variant:} values: $(p-1)! \equiv -1$ (Wilson);
$q_i! \equiv t_i! \equiv -1$ by hypothesis; $(p-1-q_i)! \equiv (-1)^{-1} =
-1$ by Lemma~\ref{lem:reflection}(i). Distinctness: cases [d], [e],
[h]--[j] as above;
$q_i = p-1$ or $t_i = p-1$ would make $p = \text{seed} + 1$ even:
impossible; $p-1-q_i = p-1 \iff q_i = 0$. All cases close, giving
$2j+e+1$ distinct points and $k = 2j+e$ intervals.

Realization of every $k$: $k = 2j+e+1$ with $j \ge 1$, $e \in \{0, 1\}$
covers all $k \ge 3$; the $-1$ variant covers $k \ge 2$. A witness for $k$
also witnesses every $k' < k$ (take a subset of the points).
\end{proof}

Note that Theorem~\ref{thm:shape} makes no claim that primes satisfying
several simultaneous seed divisibilities \emph{exist} --- that is precisely
the hard part; see Section~\ref{sec:limits}. What it shows is that
\emph{whenever} they exist, the resulting configuration is automatically
nondegenerate.

\subsection{Hexad infinitude reduced to a divisibility statement}

\begin{corollary}[Hexad reduction]\label{cor:hexad}
If there are infinitely many pairs $q_1 < q_2$ with
$q_1 \equiv q_2 \equiv 1 \pmod 6$, $q_1 \ge 7$, and
$\gcd(q_1! - 1,\ q_2! - 1) > 1$, then infinitely many primes admit hexads at
value $1$, i.e.\ Erd\H{o}s \#1056 holds for $k = 5$ (hence also $k = 3, 4$)
with infinitely many prime witnesses.

No size or congruence condition on the shared prime is needed: every prime
factor of $\gcd(q_1!-1, q_2!-1)$ automatically exceeds $q_2$ and
automatically avoids all degenerate collisions.
\end{corollary}

\begin{proof}
Let $(q_1, q_2)$ be such a pair and $p$ any prime factor of the gcd. By
Theorem~\ref{thm:shape} with $(j, e) = (2, 0)$, $p > q_2$ and $p$ carries the
hexad $\{1, q_1, q_2, p-1-q_2, p-1-q_1, p-2\}$. Given infinitely many pairs:
for fixed $q_2$ there are fewer than $q_2$ pairs, so $q_2$ takes arbitrarily
large values along the family; each pair's prime satisfies $p > q_2$, so the
set of primes produced is unbounded, hence infinite, and each carries
$k = 5$.
\end{proof}

The hypothesis is non-vacuous: nontrivial factorial gcd pairs within the
theorem-exact family exist in abundance. Verified witnesses
$(q_1, q_2;\, p)$ with $q_1 \equiv q_2 \equiv 1 \pmod 6$ and
$p \mid \gcd(q_1!-1, q_2!-1)$ include
\[
  (49,\ 97;\ 823), \quad
  (13,\ 361;\ 1733), \quad
  (415,\ 841;\ 2287), \quad
  (193,\ 265;\ 2113),
\]
\[
  (361,\ 505;\ 971), \quad
  (445,\ 1045;\ 5477);
\]
on the $-1$ side, $p = 71$ divides $q! + 1$ simultaneously for the seeds
$q = 7, 19, 61$ (all $\equiv 1 \bmod 6$; indeed $7! + 1 = 71^2$), whose
$-1$-variant configuration realizes $S_{-1}(71) = \{7, 9, 19, 51, 61, 63,
70\}$, a heptad at $-1$. Note also that $p = 971$ appears both as a pentad
prime (seed $361$, Table~\ref{tab:pentads}) and as a hexad prime (seeds
$361, 505$): the mechanisms stack. Three simultaneous seeds also occur:
$p = 4093$ divides $q! - 1$ for $q = 871, 3517, 3535$ (all
$\equiv 1 \bmod 6$) and also for the even seed $t = 2312$, a
$(j, e) = (3, 1)$ instance of Theorem~\ref{thm:shape}. It realizes
\[
  S_1(4093) = \{1,\ 557,\ 575,\ 871,\ 2312,\ 3221,\ 3517,\ 3535,\ 4091\},
\]
exactly the predicted $2j + e + 2 = 9$ points --- the three seeds, their
three mirrors $p-1-q_i$, the even seed, and the universal pair
$\{1, p-2\}$ --- giving $k = 8$ consecutive intervals at a single prime. A heuristic count (independent-values
model, corrected for the Wilson seed $q = p-2$ and reflection doubling)
predicts that the number of such pairs with $q_2 \le Q$ grows like
$Q/\log Q$ up to constants, and measured counts up to $Q = 1500$ grow
near-linearly as predicted; but no proof technique for the infinitude of
such pairs is currently visible.

\subsection{Pentads without a class number}

Theorem~\ref{thm:main} reaches the pentad through the center, and therefore
pays for it with Hypothesis~\ref{hyp:cfs}. The shape theorem offers a second,
logically independent route to the same $k = 4$, at the price of a second
divisibility instead of a class number. It is the $(j, e) = (1, 1)$ case,
which we record explicitly since it is the cheapest unrecorded corollary of
Theorem~\ref{thm:shape}.

\begin{corollary}[Mixed reduction]\label{cor:mixed}
Suppose there are infinitely many pairs $(q, t)$ with $q \equiv 1 \pmod 6$,
$q \ge 7$, with $t$ even, $t \ge 4$, and
$\gcd(q! - 1,\ t! - 1) > 1$. Then infinitely many primes admit pentads at
value $1$; that is, Erd\H{o}s \#1056 holds for $k = 4$ with infinitely many
prime witnesses. No class number condition and no congruence on $p$ is
required.
\end{corollary}

\begin{proof}
Let $(q, t)$ be such a pair and $p$ a prime factor of the gcd. By
Theorem~\ref{thm:shape} with $(j, e) = (1, 1)$, $p$ exceeds both seeds and
carries the pentad $\{1,\ q,\ t,\ p-1-q,\ p-2\}$ (sorted), all of whose
factorials are $\equiv 1 \pmod p$. For the infinitude: for a fixed bound $M$
there are fewer than $M^2$ pairs with $\max(q, t) \le M$, so $\max(q, t)$ is
unbounded along the family, and each pair's prime exceeds its seeds; the set
of primes produced is therefore unbounded, hence infinite.
\end{proof}

Two remarks. First, Corollary~\ref{cor:mixed} removes the single point of
failure in the conditional half of Theorem~\ref{thm:main}: pentad infinitude
no longer rests solely on Hypothesis~\ref{hyp:cfs}, which
Section~\ref{sec:limits} shows admits no congruence handle. Second, the
mixed hypothesis is empirically the most abundant of the family. Among the
$2262$ primes $p < 20000$, exactly $121$ carry a usable mixed pair --- an
odd seed $\equiv 1 \pmod 6$ together with an even seed $\ge 4$ --- against
only $30$ carrying two odd seeds $\equiv 1 \pmod 6$, the configuration
required by Corollary~\ref{cor:hexad}. The smallest witness is $p = 149$
with $q = 25$ and $t = 98$, giving the pentad
$\{1,\ 25,\ 98,\ 123,\ 147\}$; note $149 \equiv 1 \pmod 4$, so no central
factorial and no class number enters. This does not make the hypothesis
provable --- the obstructions of Section~\ref{sec:limits} apply to it in
the same way --- but among equally unproven hypotheses it is the one with
the widest numerical support.

\section{Limits of the identity method}\label{sec:limits}

The results above raise an obvious question: the tetrad theorem is
unconditional --- why should the pentad cost a class number condition, and
the hexad a simultaneous-divisibility hypothesis? This section proves that
the barrier is structural, not an artifact of our construction. The
``identity method'' below means: producing members of $S_1(p)$ for
infinitely many $p$ using the universal identities of
Lemma~\ref{lem:reflection} together with primes constructed as divisors of
integer sequences (Euclid-style seeds $p \mid q! \mp 1$).

\subsection{Template rigidity}

\begin{proposition}[Template rigidity]\label{prop:rigidity}
Call $\nu$ a \emph{template position} if $\nu(p) = a$ (fixed $a \ge 0$), or
$\nu(p) = p - j$ (fixed $j \ge 1$), or $\nu(p) = m \pm t$ (fixed $t \ge 0$).
Then $\nu(p)! \equiv 1 \pmod p$ holds for infinitely many primes $p$
\emph{only} in the cases
\[
  a \in \{0, 1\}, \qquad j = 2, \qquad
  (t = 0 \ \text{with} \ p \equiv 3 \bmod 4 \ \text{and} \
  h(-p) \equiv 3 \bmod 4).
\]
In every other case the set of admissible $p$ is finite and explicitly
contained in the divisors of a fixed nonzero integer:
\[
  a \ge 2:\ p \mid a! - 1; \qquad
  j \notin \{1, 2\},\ p > j:\ p \mid (j-1)! - (-1)^j;
\]
\[
  t \ge 1,\ p \equiv 3 \ (\mathrm{mod}\ 4):\ p \mid \bigl((2t-1)!!\bigr)^2 -
  4^t; \qquad
  t \ge 1,\ p \equiv 1 \ (\mathrm{mod}\ 4):\ p \mid \bigl((2t-1)!!\bigr)^2 +
  4^t.
\]
\end{proposition}

\begin{proof}[Proof sketch]
Fixed low positions $a \ge 2$ are immediate. Top positions: by
Lemma~\ref{lem:reflection}(iv), $(p-j)! \equiv 1 \iff p \mid (j-1)! -
(-1)^j$, a nonzero integer for $j \notin \{1, 2\}$ since $(j-1)! \ge 2$
($j = 1$ gives $p \mid 2$; $j = 2$ gives $0$, i.e.\ all $p$). Near-center
positions: the telescoped central formulas
\[
  (m+t)! \equiv m! \cdot \frac{(2t-1)!!}{2^t}, \qquad
  (m-t)! \equiv m! \cdot \frac{(-2)^t}{(2t-1)!!} \pmod p
\]
(both consequences of $j \equiv -(p - j)$ pairing around the center) combine
with $(m!)^2 \equiv (-1)^{m+1}$: for $p \equiv 3 \pmod 4$,
$m! = \sigma \in \{\pm 1\}$ and $(m \pm t)! \equiv 1$ forces
$p \mid (2t-1)!! \mp \sigma(\pm 2)^t$, hence in either case
$p \mid ((2t-1)!!)^2 - 4^t$, nonzero (odd square minus even power) for
$t \ge 1$; for $p \equiv 1 \pmod 4$, $(m!)^2 \equiv -1$ and squaring gives
$p \mid ((2t-1)!!)^2 + 4^t$. The center itself ($t = 0$): for
$p \equiv 1 \pmod 4$, $(m!)^2 \equiv -1$, so $m!$ has order $4$ and is never
$\pm 1$; for $p \equiv 3 \pmod 4$ the sign is $h(-p) \bmod 4$
(Lemmas~\ref{lem:centralsign}--\ref{lem:bridge}), which admits no known
congruence description: empirically both signs occur with frequency
$\tfrac12 \pm 0.03$ in every residue class modulo $8$, $16$, and $24$ for
$p \le 60000$.
\end{proof}

\subsection{Bounded patterns are finite}

A \emph{pattern} is a composition $(L_1, \dots, L_t)$ of a total length
$C = L_1 + \dots + L_t$, describing $t$ adjacent intervals of consecutive
integers of lengths $L_i$, placed at a variable base: with cut points
$c_0 = 0 < c_1 < \dots < c_t = C$ and base $m$, the chain condition
$(m + c_0)! \equiv (m + c_1)! \equiv \dots \equiv (m + c_t)! \pmod p$
(for $p > m + C$) is equivalent to the polynomial system
\[
  g_i(m) \equiv 0 \pmod p, \qquad
  g_i(x) = \prod_{j = c_{i-1}+1}^{c_i} (x + j) \; - \; 1, \qquad
  i = 1, \dots, t .
\]

\begin{theorem}[Bounded-pattern finiteness; computational]\label{thm:patterns}
For every one of the $65{,}519$ patterns with $C \le 16$ and $t \ge 2$, and
every one of the $1{,}966{,}080$ resulting polynomial pairs $(g_i, g_j)$,
$\gcd(g_i, g_j) = 1$ in $\Q[x]$. Consequently, for every such pattern the
primes admitting a chain with all cut points inside a window of length
$\le 16$ divide the (nonzero) pairwise resultants $\Res(g_i, g_j)$ --- a
finite, explicit set per pattern. There is no bounded-pattern route to
infinitude for any $k \ge 2$, and no ``algebraic jackpot'' pattern exists at
small size.
\end{theorem}

\begin{proof}[Proof sketch]
The equivalence chain-condition $\leftrightarrow$ polynomial system follows
from $(m + c_i)! = (m + c_{i-1})! \prod_{j = c_{i-1}+1}^{c_i} (m + j)$ and
$p \nmid (m + c_{i-1})!$. For the coprimality certificate: reducing modulo a
prime $\ell$ can only increase the degree of a gcd (Sylvester/subresultant
argument), so exhibiting $\deg \gcd_{\F_\ell}(g_i, g_j) = 0$ for a single
$\ell$ proves coprimality over $\Q$; the computation used
$\ell = 2^{61} - 1$ for every pair. If $\gcd(g_i, g_j) = 1$ then
$\Res(g_i, g_j) \ne 0$, and any prime modulo which $g_i, g_j$ share a root
divides it.
\end{proof}

The historically known sporadic witnesses are precisely such resultant
primes: for the pattern $(2, 3)$, $\Res(g_1, g_2) = 11$ exactly ---
Erd\H{o}s's example $3 \cdot 4 \equiv 5 \cdot 6 \cdot 7 \equiv 1 \pmod{11}$
lives at the root $m \equiv 2 \pmod{11}$ --- and $17$ divides all three
pairwise resultants of Makowski's pattern $(4, 6, 4)$, root
$m \equiv 1 \pmod{17}$. In this light, A060427 is a record sequence drawn
from prime factors of resultant systems of provably coprime polynomial
families, plus unstructured coincidences: sporadic by nature. For $C > 16$
pairwise coprimality remains conjectural; it is a close cousin of the
conjecture of Klurman and Munsch \cite{KM} (Conjecture 3.3 in the published
version; Conjecture 4.2 in arXiv:1505.01198) that the
splitting fields of $t(t+1)\cdots(t+n-1) - 1$ for distinct $n$ intersect
only in $\Q$ --- that is, the experts' conjecture is that this obstruction is
\emph{permanent}.

\subsection{Cheap second divisibilities do not exist}

\begin{lemma}\label{lem:gcds}
For all integers $q, n \ge 2$:
\[
  \gcd(q!-1,\ (q+1)!-1) = 1; \qquad
  \gcd(q!-1,\ q!+1) = 1; \qquad
  \gcd(n!+1,\ (n+1)!+1) = 1;
\]
and $\gcd(q!-1, (q+1)!+1)$ divides $q+2$, hence equals $1$ for
$q \equiv 1 \pmod 6$, $q \ge 7$.
\end{lemma}

\begin{proof}
If $d \mid q!-1$ and $d \mid (q+1)!-1$ then
$d \mid (q+1)(q!-1) - ((q+1)!-1) = -q$, so $d \mid q$; but $d \mid q!-1$
forces $\gcd(d, q!) = 1$, so $d = 1$. The second gcd divides $2$ and both entries
are odd. The third: $d \mid (n+1)(n!+1) - ((n+1)!+1) = n$, and
$\gcd(d, n!) = 1$ again. The fourth:
$d \mid (q+1)!+1 - (q+1)(q!-1) = q+2$; for $q \equiv 1 \pmod 6$ every prime
factor of $q+2 = 3 \cdot \frac{q+2}{3}$ is $\le \frac{q+2}{3} < q$ or equals
$3$, while every prime factor of $q!-1$ exceeds $q$; so $d = 1$.
\end{proof}

\subsection{The bootstrap ceiling}

\begin{proposition}[Bootstrap]\label{prop:bootstrap}
Fix $t \ge 1$ and odd seeds $q_1 < \dots < q_t$. The $t$ simultaneous
conditions $q_i! \equiv 1 \pmod p$ are equivalent to
\[
  p \mid q_1! - 1
  \quad\text{together with}\quad
  p \ \Bigm|\ \prod_{q_{i-1} < j \le q_i} j \; - \; 1, \qquad
  i = 2, \dots, t
\]
--- that is, to a witness of the \emph{anchored $t$-interval instance of
Erd\H{o}s \#1056 itself}. Hence the identity method converts infinitude for
the $t$-seed configuration ($k = 2t + 1$ intervals by
Theorem~\ref{thm:shape}) into infinitude for a $t$-interval instance of the
same problem, and cannot self-start beyond $t = 1$: the case $t = 1$ is
exactly the single-sequence Euclid construction (the tetrad theorem of
\cite{AAHS}, $k = 3$), while every $t \ge 2$ presupposes an instance of the
problem being solved.
\end{proposition}

\begin{proof}
$p \mid q_1! - 1$ implies $p \nmid q_i!$ for all $i$, so
$q_i! \equiv 1 \iff$ the interval product over $(q_{i-1}, q_i]$ is
$\equiv 1$; conversely those interval conditions are precisely ``$t-1$
further consecutive intervals of product $\equiv 1$, anchored at
$q_1 \in S_1(p)$''.
\end{proof}

\subsection{The exhaustion theorem and the equivalence}

\begin{theorem}[Exhaustion]\label{thm:exhaustion}
Fix $W \ge 3$ and let $C(W)$ be the maximum of the finite set
\[
  \{n! - 1 : 2 \le n \le W\} \cup \{|(j-1)! - (-1)^j| : 3 \le j \le W\}
  \cup \{\,\bigl|((2t-1)!!)^2 \mp 4^t\bigr| : 1 \le t \le W\,\}.
\]
Then for every prime $p > C(W)$, with $m = (p-1)/2$:
\[
  S_1(p) \,\cap\, \Bigl([1, W] \cup [p-W,\, p-1] \cup [m-W,\, m+W]\Bigr)
  \;=\;
  \{1,\ p-2\} \ \cup\
  \begin{cases}
    \{m\} & \text{if } p \equiv 3 \ (\mathrm{mod}\ 4) \text{ and }
             h(-p) \equiv 3 \ (\mathrm{mod}\ 4), \\
    \emptyset & \text{otherwise.}
  \end{cases}
\]
\end{theorem}

\begin{proof}[Proof sketch]
Each window is handled by Proposition~\ref{prop:rigidity}: a member of
$S_1(p)$ at position $n \in [2, W]$, $n = p-j$ with $3 \le j \le W$, or
$n = m \pm t$ with $1 \le t \le W$ forces $p$ to divide a nonzero integer
bounded by $C(W) < p$, which is impossible; the surviving positions are
exactly $1$, $p-2$, and the conditional center.
\end{proof}

Assembling Proposition~\ref{prop:rigidity},
Theorem~\ref{thm:patterns}, Lemma~\ref{lem:gcds},
Proposition~\ref{prop:bootstrap}, and Theorem~\ref{thm:exhaustion}:

\begin{theorem}[Limitative theorem]\label{thm:limitative}
Within the identity framework --- members of $S_1(p)$ produced by the
universal identities of Lemma~\ref{lem:reflection} at template positions,
with primes obtained as divisors of constructed integers --- the following
hold.
\begin{enumerate}
  \item[(i)] One divisibility $p \mid q!-1$ ($q$ odd) yields exactly the
    orbit $\{q, p-1-q\}$ of the Wilson involution, on top of the universal
    pair $\{1, p-2\}$: four elements, the tetrad, $k = 3$. This is the
    unconditional ceiling of the method.
  \item[(ii)] A fifth element \emph{at value $1$} requires the center, i.e.\
    the open statement $h(-p) \equiv 3 \pmod 4$ infinitely often along
    $\mathcal F$ (Hypothesis~\ref{hyp:cfs}); there is no congruence
    handle on this sign. (At a free common value the center is available
    unconditionally; see Remark~\ref{rem:freevalue}.)
  \item[(iii)] The sixth element requires a second independent divisibility
    $p \mid \gcd(q_1!-1, q_2!-1)$; cheap adjacent variants are provably
    trivial (Lemma~\ref{lem:gcds}), bounded patterns admit only finitely
    many primes each (Theorem~\ref{thm:patterns}), and unbounded second
    seeds re-pose the problem itself (Proposition~\ref{prop:bootstrap}).
\end{enumerate}
Hence, \emph{at the pinned value $1$}, pentad or hexad infinitude via
identities is \emph{equivalent} to the class number statement
(Hypothesis~\ref{hyp:cfs}) or to a second seed divisibility (hypothesis of
Corollary~\ref{cor:mixed} or~\ref{cor:hexad}), respectively. This explains
precisely why the June 2026 state of the art ($k = 3$) is the natural
ceiling of the classical toolkit \emph{as applied at a pinned value}.
\end{theorem}

We regard Theorem~\ref{thm:limitative} as the real content of this note
alongside Theorem~\ref{thm:main}: not a proof that $k \ge 4$ infinitude is
false --- we believe the opposite --- but an identification of the exact
open statements any identity-based proof must pass through.

\begin{remark}[The scope of the framework: free common values]
\label{rem:freevalue}
The restriction to $S_1(p)$ in the definition of the identity method is
essential and not cosmetic, and we emphasize it because Erd\H{o}s \#1056
does \emph{not} pin the common value: it asks only for consecutive intervals
of product $\equiv 1$, which is the statement that $k+1$ points share
\emph{some} common factorial residue. The near-center identities displayed
in the proof of Proposition~\ref{prop:rigidity} are themselves
value-agnostic. Reading them without imposing $(m \pm t)! \equiv 1$, one
gets at once, for \emph{even} $t$ with $1 \le t \le m-1$,
\[
  (2t-1)!! \equiv 2^t \pmod p
  \quad\Longrightarrow\quad
  (m-t)! \equiv m! \equiv (m+t)! \pmod p,
\]
since the two displayed ratios become $1$ and $(-1)^t = 1$ respectively.
This costs no congruence on $p$ and no class number: the shared residue is
whatever $m!$ happens to be. Over the primes $p < 20000$ and even
$t \le 400$ there are $142$ such $(p, t)$, all confirmed by direct
computation, and in $46$ of them the shared value is neither $+1$ nor $-1$
--- configurations lying entirely outside $S_1(p) \cup S_{-1}(p)$.

The consequence for the present section is a limitation of scope rather than
an error: Theorem~\ref{thm:limitative} bounds what the identity method
achieves \emph{at a pinned value}, and says nothing about free-value
configurations. These can be large. At $p = 977$ we find
\[
  102!\ \equiv\ 330!\ \equiv\ 483!\ \equiv\ 488!\ \equiv\ 646!\ \equiv\
  874!\ \equiv\ 725 \pmod{977},
\]
a hexad ($k = 5$) at a prime $\equiv 1 \pmod 4$, at a common value that is
neither $\pm 1$. Here $488 = m$, the pairs $(330, 646)$ and $(102, 874)$ are
$m \mp t$ for $t = 158$ and $t = 386$ --- both instances of the implication
above --- while $483 = m - 5$ is an additional coincidence with no mechanism
we can attach to it. Such configurations are the norm rather than the
exception: of the $2259$ primes in $[7, 20000)$, exactly $2207$ --- some
$97.7\%$ --- carry $k \ge 4$ at some common value, while only $344$
($15.2\%$) carry $k \ge 4$ at the value $1$. The wall at $k = 3$ is thus a
wall of \emph{provability}, not of scarcity, and the free-value regime is
where we would now look for a way around it.
\end{remark}

\section{Numerical evidence}\label{sec:numerics}

\subsection{Certified pentads}

Table~\ref{tab:pentads} lists twelve pentads certified end to end: the seed
divisibility $p \mid q!-1$, the primality of $p$, the congruences
$p \equiv 3 \pmod 4$ and $q \equiv 1 \pmod 6$, the five factorial values,
distinctness, and the class number $h(-p) \equiv 3 \pmod 4$ were each
verified by direct computation, with $h(-p)$ computed two independent ways
(exact count of reduced primitive binary quadratic forms of discriminant
$-p$, and the Dirichlet character sum), in agreement in every case. Note
that composite seeds are allowed and occur ($q = 25 = 5^2$, $q = 115$,
$q = 235$, $q = 295$, $q = 301$, $q = 361 = 19^2$).

\begin{table}[ht]
\centering
\begin{tabular}{rrrrrr}
\hline
$p$ & $q$ & $m = \frac{p-1}{2}$ & $p-1-q$ & $h(-p)$ & \\
\hline
$359$   & $229$ & $179$   & $129$   & $19$  & \\
$439$   & $37$  & $219$   & $401$   & $15$  & \\
$883$   & $43$  & $441$   & $839$   & $3$   & \\
$907$   & $25$  & $453$   & $881$   & $3$   & \\
$971$   & $361$ & $485$   & $609$   & $15$  & \\
$5003$  & $295$ & $2501$  & $4707$  & $15$  & \\
$5039$  & $7$   & $2519$  & $5031$  & $83$  & (flagship: $5039 = 7!-1$) \\
$9439$  & $235$ & $4719$  & $9203$  & $75$  & \\
$9643$  & $337$ & $4821$  & $9305$  & $11$  & \\
$9791$  & $301$ & $4895$  & $9489$  & $119$ & \\
$21839$ & $115$ & $10919$ & $21723$ & $195$ & \\
$22639$ & $223$ & $11319$ & $22415$ & $55$  & \\
\hline
\end{tabular}
\caption{Twelve certified pentads. In each row the pentad is
$\{1, q, m, p-1-q, p-2\}$, all five factorials $\equiv 1 \pmod p$, and
$h(-p) \equiv 3 \pmod 4$ as the criterion requires.}
\label{tab:pentads}
\end{table}

\subsection{The flagship}

The prime $p = 5039 = 7! - 1$ is itself of the form $q! - 1$ with the seed
$q = 7 \equiv 1 \pmod 6$, and $5039 \equiv 3 \pmod 4$ with
$h(-5039) = 83 \equiv 3 \pmod 4$. An exhaustive scan of $n = 1, \dots,
5038$ shows that the complete solution set of $n! \equiv 1 \pmod{5039}$ is
\emph{exactly} the constructed pentad
\[
  \{\, 1,\ 7,\ 2519,\ 5031,\ 5037 \,\},
\]
and $n! \equiv -1 \pmod{5039}$ only at the Wilson point $n = 5038$. Thus for
this prime the theorem does not merely produce five solutions; it accounts
for every solution.

\subsection{Tightness in the negative direction}

The criterion bites both ways. For $q = 31$, the divisor $p = 787$ of
$31! - 1$ has $h(-787) \equiv 1 \pmod 4$; for $q = 13$, the only prime
factor $\equiv 3 \pmod 4$ of $13! - 1$ is $p = 3593203$, with
$h(-p) \equiv 1 \pmod 4$. In both cases $m! \equiv -1 \pmod p$ was verified
directly, and the scan confirmed that no pentad through $m$ exists: case
(ii) of Theorem~\ref{thm:main}, exactly as predicted.

\subsection{Sign statistics}

Across all $1632$ primes $p \equiv 3 \pmod 4$ in $[7, 30000]$, the central
sign $m! \equiv \pm 1 \pmod p$ splits $826$ at $+1$ to $806$ at $-1$
--- consistent with the expected equidistribution of $h(-p) \bmod 4$
between $1$ and $3$. The split remains balanced (within $\pm 0.03$ of
$\tfrac12$) when refined by residue classes modulo $8$, $16$, and $24$ for
$p \le 60000$; no congruence bias is visible, which is why
Hypothesis~\ref{hyp:cfs} appears both true and inaccessible to congruence
methods.

\subsection{Independent verification}

The elementary components were stress-tested independently of the proofs:
the reflection identity of Lemma~\ref{lem:reflection} was verified
exhaustively for all $n \in [0, p-1]$ for every prime $p < 500$ and for
$p = 10007$ and $p = 99991$; the Mordell law
(Lemmas~\ref{lem:centralsign}--\ref{lem:bridge}) was verified for all
$1135$ primes $p \equiv 3 \pmod 4$ in $[7, 20000)$, with class numbers
computed by reduced-form counting --- millions of individual assertions in
total, with zero failures.

\begin{remark}[Provenance, verification status, and honesty]
\label{rem:provenance}
This note reports work carried out in an AI-assisted setting: the proof
search was conducted by multiple independent prover agents, the results were
attacked by adversarial refuter agents, and all computations described above
were machine-executed and machine-checked. The reader should weigh the
following plainly stated caveats. (1) The elementary core
(Lemmas~\ref{lem:reflection}--\ref{lem:bridge} and
Theorem~\ref{thm:main}) is short enough to verify by hand, and a Lean
formalization of this core is in progress but not complete at the time of
writing. (2) The novelty assessment --- that the coupling of Erd\H{o}s
\#1056 to $h(-p) \bmod 4$ does not appear in the tetrads development
\cite{AAHS}, in the discussion at erdosproblems.com/1056 \cite{EP}, or in
Hardy--Subbarao \cite{HS} --- rests on machine literature searches and a
dedicated follow-up check, not on a systematic human review; the coupling
is natural and prior art may exist. (3) The Wilson pairing toolkit itself is
classical and no novelty is claimed for it, nor for Mordell's theorem, of
which Lemmas~\ref{lem:centralsign} and \ref{lem:bridge} are a re-proof.
(4) The tetrad construction and the family $q \equiv 1 \pmod 6$,
$p \mid q!-1$ are entirely due to Agust\'in-Aquino and Hern\'andez
Santiago \cite{AAHS}; our theorem adds the central point and the class
number dichotomy. (5) Above all: \emph{no unconditional infinitude for any
$k \ge 4$ is claimed}. Corollary~\ref{cor:cond} is conditional on
Hypothesis~\ref{hyp:cfs}, Corollary~\ref{cor:hexad} on a divisibility
hypothesis, and Theorem~\ref{thm:limitative} explains why some such
hypothesis is currently unavoidable. The Erd\H{o}s problem remains open.
\end{remark}

\section*{Acknowledgments}
The tetrad theorem and its Lean development by O.~A.~Agust\'in-Aquino and
J.~Hern\'andez Santiago \cite{AAHS} are the foundation on which this note
builds. Thanks are due to the maintainers of erdosproblems.com and the OEIS,
whose curation makes attacks like this one possible.

\begin{thebibliography}{9}

\bibitem{AAHS}
O.~A.~Agust\'in-Aquino and J.~Hern\'andez Santiago,
\emph{Infinite tetrads of congruent factorials} (Lean-verified development
proving the infinitude of tetrad primes),
2026. Repository: \url{https://github.com/octavioalberto/tetrads}.

\bibitem{BS}
Z.~I.~Borevich and I.~R.~Shafarevich,
\emph{Number Theory},
Academic Press, New York, 1966. (Counting form of Dirichlet's class number
formula for imaginary quadratic fields.)

\bibitem{EP}
T.~F.~Bloom (maintainer),
\emph{Erd\H{o}s problem \#1056},
\url{https://www.erdosproblems.com/1056}.

\bibitem{Guy}
R.~K.~Guy,
\emph{Unsolved Problems in Number Theory},
3rd ed., Problem A15, Springer-Verlag, New York, 2004.

\bibitem{HS}
G.~E.~Hardy and M.~V.~Subbarao,
\emph{A modified problem of Pillai and some related questions},
Amer. Math. Monthly \textbf{109} (2002), 554--559.

\bibitem{KM}
O.~Klurman and M.~Munsch,
\emph{Distribution of factorials modulo $p$},
J. Th\'eor. Nombres Bordeaux \textbf{29} (2017), 169--177.
arXiv:1505.01198.

\bibitem{Landau}
E.~Landau,
\emph{Vorlesungen \"uber Zahlentheorie},
Hirzel, Leipzig, 1927. (Classical counting form of the class number
formula.)

\bibitem{Mordell}
L.~J.~Mordell,
\emph{The congruence $((p-1)/2)! \equiv \pm 1 \pmod p$},
Amer. Math. Monthly \textbf{68} (1961), 145--146.

\bibitem{OEIS}
OEIS Foundation Inc.,
\emph{Sequence A060427: smallest prime $p$ such that there are $n$ strings
of consecutive integers all having products $\equiv 1 \bmod p$},
The On-Line Encyclopedia of Integer Sequences,
\url{https://oeis.org/A060427}.

\bibitem{Wi79}
K.~S.~Williams,
\emph{The class number of $\mathbb{Q}(\sqrt{-p})$ modulo $4$, for
$p \equiv 3 \pmod 4$ a prime},
Pacific J. Math. \textbf{83} (1979), no.~2, 565--570.

\end{thebibliography}

\end{document}
